\section{Experiments}
\label{sec:experiments}

\subsection{Experimental Setup}
\label{sec:experimental-setup}

We conduct experiments under the same rationale-based distillation setting used in the original study. In this setting, rationales generated by a teacher LLM are used as auxiliary supervision for training a compact student model. The student is trained to predict the task label and to generate the rationale, while inference requires only the student model. This paper keeps that training framework unchanged and evaluates whether the proposed boundary-aware selection stage improves the quality of the rationales used for supervision.

Following the original comparison, we evaluate on ESNLI and CommonsenseQA. ESNLI evaluates textual entailment over premise-hypothesis pairs, and CommonsenseQA evaluates multiple-choice commonsense reasoning. Accuracy is used as the evaluation metric for both tasks.

\begin{table}[H]
\centering
\caption{Datasets used in the experiments.}
\label{tab:experiment-datasets}
\begin{tabular}{p{0.18\linewidth} p{0.28\linewidth} p{0.24\linewidth} p{0.22\linewidth}}
\hline
Dataset & Task & Output & Role in this paper \\
\hline
ESNLI & Natural language inference & entailment / neutral / contradiction & Main textual entailment evaluation \\
\hline
CommonsenseQA & Commonsense question answering & multiple-choice answer & Additional reasoning/generalization evaluation \\
\hline
\end{tabular}
\end{table}

The baseline is the original rationale-distilled student without boundary-aware selection. The main proposed method is Boundary-Mix, which keeps one selected rationale from easy and boundary examples and removes hard examples. Boundary-Bridge and Boundary-Specialist are evaluated as same-family variants to isolate the effect of changing the final selection rule while keeping the shared boundary pipeline fixed.

\subsection{Overall Performance Compared with Previous Models}
\label{sec:overall-performance}

We first compare our method with the reference models reported in the original experiment. This comparison shows whether boundary-aware rationale selection preserves the main advantage of rationale-based distillation: a compact student model can outperform larger or more expensive baselines when it receives useful reasoning supervision. The most direct comparison is between the original rationale-distilled student and our method, because both use the same compact-student setting and differ mainly in how training rationales are selected.

\begin{table}[H]
\centering
\caption{Overall model comparison. Accuracy is reported in percentage.}
\label{tab:overall-performance}
\begin{tabular}{p{0.34\linewidth} r r r p{0.18\linewidth}}
\hline
Model / Method & Parameters & ESNLI & CommonsenseQA & Source \\
\hline
RoBERTa Base & 125M & 44.08 & -- & Original comparison \\
FLAN-T5 Large & 783M & 57.29 & 60.11 & Original comparison \\
Standard Finetune T5 Base & 220M & 61.41 & 52.74 & Original comparison \\
ChatGPT 3.5 & 175B & 72.99 & 77.23 & Original comparison \\
Distilling Step-by-Step & 220M & 75.38 & 55.45 & Original comparison \\
Original rationale-distilled student & 220M & 83.93 & 60.36 & Current baseline \\
\textbf{Our method} & 220M & \textbf{84.16} & \textbf{62.41} & Boundary-aware selection \\
\hline
\end{tabular}
\end{table}

Our method achieves \textbf{84.16\%} accuracy on ESNLI and \textbf{62.41\%} on CommonsenseQA. Compared with the original rationale-distilled student, it improves ESNLI from 83.93\% to 84.16\% and CommonsenseQA from 60.36\% to 62.41\%. It also remains far smaller than ChatGPT 3.5 while outperforming it on ESNLI in this comparison, and it improves over FLAN-T5 Large on both datasets despite using a smaller student model. These results support the main claim that selecting more suitable rationales can strengthen rationale-based distillation without changing the student architecture.

The broad model table gives context. The controlled comparison is the improvement from the original rationale-distilled student to our method.

\subsection{Boundary Family Performance}
\label{sec:boundary-family-performance}

We next compare variants from the same boundary-aware family. This comparison is important because the variants are not unrelated methods: they share candidate grouping, weighted voting, confidence estimation, difficulty-band assignment, and rationale scoring. Their difference is the final selection rule. Therefore, changes in performance can be interpreted as the effect of keeping one primary rationale, adding a bridge rationale, or focusing on specialist sources.

\begin{table}[H]
\centering
\caption{Boundary-family performance. Accuracy is reported in percentage.}
\label{tab:boundary-family-performance}
\begin{tabular}{p{0.20\linewidth} p{0.34\linewidth} r r p{0.22\linewidth}}
\hline
Method & Final selection behavior & ESNLI & CommonsenseQA & Interpretation \\
\hline
Baseline & No boundary-aware filtering & 83.93 & 60.36 & Original rationale-distilled student \\
\hline
Boundary-Mix & Keep easy + boundary; one rationale per example & 84.16 & 62.41 & Main method; improves both datasets \\
\hline
Boundary-Bridge & Keep easy + boundary and add bridge rationale when useful & 85.06 & 60.52 & Best ESNLI, weak CommonsenseQA gain \\
\hline
Boundary-Specialist & Focus on boundary/bridge rows from specialist sources & 83.60 & training in progress & ESNLI-only result so far; useful as source-focused analysis \\
\hline
\end{tabular}
\end{table}

Boundary-Bridge gives the highest ESNLI score, which suggests that sentence-pair entailment can benefit from a second complementary rationale. However, its CommonsenseQA result is only slightly above the baseline, indicating that an additional rationale may introduce less useful associations for multiple-choice commonsense reasoning. Boundary-Specialist reaches 83.60\% on ESNLI, lower than the baseline and the other two boundary variants, which suggests that focusing only on specialist/boundary-style sources removes too much clean supervision. Boundary-Mix is therefore the main method because it improves both completed datasets and gives the strongest CommonsenseQA result among the evaluated family variants.

The selection statistics support this interpretation. Boundary-Mix keeps a conservative training set with one selected rationale per retained example. Boundary-Bridge expands the selected data by adding bridge rationales, which increases coverage and helps ESNLI. Boundary-Specialist is much smaller and excludes easy examples, which explains why it is useful for source analysis but weaker as a general training set.

\subsection{Rationale-Type Distribution Across Family Variants}
\label{sec:rationale-type-distribution}

We further analyze which rationale types are selected by each family variant. This analysis explains the composition of the selected supervision rather than the final model accuracy. Tables~\ref{tab:esnli-rationale-distribution} and~\ref{tab:cqa-rationale-distribution} report the selected rationale-source distributions.

\begin{table}[H]
\centering
\caption{Selected rationale-source distribution on ESNLI. Values are percentages of selected rows.}
\label{tab:esnli-rationale-distribution}
\begin{tabular}{p{0.22\linewidth} r r r}
\hline
Rationale source & Boundary-Mix & Boundary-Bridge & Boundary-Specialist \\
\hline
causal & 38.9 & 33.6 & -- \\
neutral & 28.3 & 31.4 & -- \\
if-else & 18.2 & 16.0 & 30.7 \\
contrastive & 5.3 & 7.2 & 25.6 \\
comparative & 4.1 & 5.5 & 20.5 \\
historical & 2.7 & 3.0 & 10.0 \\
consensus & 2.6 & 3.3 & 13.3 \\
\hline
\end{tabular}
\end{table}

\begin{table}[H]
\centering
\caption{Selected rationale-source distribution on CommonsenseQA. Values are percentages of selected rows.}
\label{tab:cqa-rationale-distribution}
\begin{tabular}{p{0.24\linewidth} r r r}
\hline
Rationale source & Boundary-Mix & Boundary-Bridge & Boundary-Specialist \\
\hline
if-else & 77.9 & 30.7 & 48.7 \\
causal & 12.1 & 30.7 & 0.0 \\
neutral & 8.7 & 30.7 & 0.0 \\
contrastive & 1.3 & 7.7 & 49.3 \\
historical & 0.0 & 0.1 & 0.8 \\
consensus & 0.0 & 0.1 & 0.8 \\
comparative & -- & 0.0 & 0.3 \\
\hline
\end{tabular}
\end{table}

Tables~\ref{tab:esnli-rationale-distribution} and~\ref{tab:cqa-rationale-distribution} show two different behaviors. On ESNLI, Boundary-Mix and Boundary-Bridge both keep a broad mixture of causal, neutral, and if-else rationales, while Boundary-Specialist shifts toward if-else, contrastive, and comparative sources. On CommonsenseQA, Boundary-Mix is strongly dominated by if-else rationales, whereas Boundary-Bridge distributes selected rationales more evenly across causal, if-else, and neutral sources. Boundary-Specialist is concentrated almost entirely on contrastive and if-else rationales. This supports the family-performance interpretation: ESNLI benefits from rationale diversity, while CommonsenseQA is more stable when the method keeps a single clean rationale type.

\subsection{Selection Statistics Behind the Distributions}
\label{sec:selection-statistics}

Table~\ref{tab:selection-statistics} reports supporting selection statistics. This table explains how many rows each variant contributes and whether the selected rationales are mainly easy, boundary, or bridge examples.

\begin{table}[H]
\centering
\caption{Selection statistics for boundary-family variants.}
\label{tab:selection-statistics}
\begin{tabular}{p{0.16\linewidth} p{0.18\linewidth} r p{0.26\linewidth} r r}
\hline
Dataset & Variant & Selected rows & Band distribution & Avg. agreement & Avg. margin \\
\hline
ESNLI & Boundary-Mix & 9054 & easy: 9032, boundary: 22 & 6.85 & 5.95 \\
ESNLI & Boundary-Bridge & 12600 & easy: 7747, bridge: 4853 & 6.94 & 6.08 \\
ESNLI & Boundary-Specialist & 3921 & bridge: 3912, boundary: 9 & 6.84 & 5.93 \\
CommonsenseQA & Boundary-Mix & 7843 & easy: 7787, boundary: 56 & 6.67 & 5.98 \\
CommonsenseQA & Boundary-Bridge & 10500 & bridge: 6037, easy: 4462, boundary: 1 & 6.91 & 6.37 \\
CommonsenseQA & Boundary-Specialist & 1437 & bridge: 1402, boundary: 35 & 6.54 & 5.78 \\
\hline
\end{tabular}
\end{table}

CommonsenseQA Boundary-Specialist selection statistics are included here, although its downstream training performance is not yet reported.

\subsection{Qualitative Error Analysis}
\label{sec:qualitative-error-analysis}

Finally, we follow the qualitative style of the original paper and inspect concrete prediction examples from the ESNLI test output. This analysis is important because the selected rationale can be linguistically fluent but still incomplete for the final label decision. In the ESNLI prediction file, the student reaches 84.16\% accuracy, while the raw LLM label reaches 72.99\% accuracy after normalizing labels. The comparison between the LLM label and the student prediction is summarized in Table~\ref{tab:esnli-error-counts}.

\begin{table}[H]
\centering
\caption{Agreement between raw LLM labels and student predictions on ESNLI test examples.}
\label{tab:esnli-error-counts}
\begin{tabular}{lrr}
\hline
Outcome category & Count & Interpretation \\
\hline
LLM correct, student correct & 6281 & Both teacher signal and student prediction agree with the gold label. \\
LLM wrong, student correct & 1987 & The student corrects or ignores a noisy teacher label. \\
LLM correct, student wrong & 890 & The teacher signal is correct, but the student fails to learn the decision. \\
LLM wrong, student wrong & 666 & Both teacher signal and student prediction fail. \\
\hline
\end{tabular}
\end{table}

Table~\ref{tab:qualitative-examples} shows representative examples from these four groups. For traceability, each row also reports the original row index from \texttt{result/esnli\_test\_predictions.csv}. For each example, we report the gold label, normalized LLM label, student prediction, rationale, and interpretation.

\begin{table}[H]
\centering
\caption{Representative ESNLI examples for qualitative error analysis.}
\label{tab:qualitative-examples}
\begin{tabular}{p{0.06\linewidth} p{0.13\linewidth} p{0.075\linewidth} p{0.075\linewidth} p{0.075\linewidth} p{0.33\linewidth} p{0.18\linewidth}}
\hline
Row & Case & Gold & LLM & Student & Input and rationale & Interpretation \\
\hline
1 &
Both correct &
entailment &
entailment &
entailment &
Premise: ``This church choir sings to the masses as they sing joyous songs from the book at a church.'' Hypothesis: ``The church is filled with song.'' Rationale: ``The church choir is singing joyous songs from the book.'' &
The student predicts entailment, which is correct because singing joyous songs in the church supports the hypothesis that the church is filled with song. \\
\hline
9355 &
LLM wrong, student correct &
contradiction &
neutral &
contradiction &
Premise: ``A dog jumping over a beam.'' Hypothesis: ``An animal is lying down.'' Rationale: ``A dog is an animal.'' &
The student predicts contradiction, which is correct because ``jumping over a beam'' conflicts with ``lying down.'' The LLM label is neutral because its rationale only captures the entity relation. \\
\hline
4076 &
LLM correct, student wrong &
entailment &
entailment &
contradiction &
Premise: ``Two daschunds play with a red ball.'' Hypothesis: ``Two dogs playing together.'' Rationale: ``Daschunds are dogs.'' &
The student predicts contradiction, which is wrong because ``daschunds are dogs'' supports the entailment relation. This shows a remaining lexical-category failure. \\
\hline
7707 &
Both wrong &
neutral &
entailment &
entailment &
Premise: ``Some youngsters are playing a card game.'' Hypothesis: ``people play poker.'' Rationale: ``youngsters are people.'' &
The student predicts entailment, which is wrong because a card game does not necessarily mean poker. The rationale supports ``youngsters are people'' but not the specific game. \\
\hline
\end{tabular}
\end{table}

These cases show why the proposed selection stage should evaluate more than rationale fluency. Some rationales contain a true statement but do not justify the full label, such as ``A dog is an animal'' for a contradiction example or ``youngsters are people'' for a neutral example. Boundary-aware selection is designed to reduce this kind of noisy supervision by checking candidate agreement, label margin, source compatibility, and rationale quality before the example is used for student training.
